% Vorlage für Abschlussarbeiten
%
% einstellungen.tex
%
% Hier werden allgemeine Einstellungen festgelegt.

\documentclass[%
	a4paper,%							A4 Papier
	oneside,%							einseitig (linker und rechter Seitenrand sind gleich groß)
	bibliography=totoc,%	Literaturverzeichnis nummeriert mit ins Inhaltsverzeichnis einfügen
	listof=totoc,%				Abbildungs- und Tabellenverzeichnis ins Inhaltsverzeichnis einfügen
	numbers=noenddot,%		hinter der Gliederungsnummer soll kein Punkt gesetzt werden (siehe Duden)
	parskip=false,%				europäischer Satz mit Abstand zwischen Absätzen
	headsepline,%					Linie nach Kopfzeile
	captions=tableheading,%		Tabellenueberschriften mit passendem Abstand setzen
	headheight=17pt,%
	chapterprefix=false,%	kein "Kapitel N" über der Kapitelüberschrift
	11pt,%								11 pt Grundschrift (Leitfäden: 11–12 pt), mit \onehalfspacing 16,5 pt Zeilenabstand = Word 1,3
]{scrreprt}

% Laut Duden gilt:
% Bei rein arabischer Nummerierung soll der Punkt
% weggelassen werden. Werden auch römische Zahlen und
% Buchstaben für die Gliederung verwendet, soll ein
% abschließender Punkt gesetzt werden. Als Voreinstellung
% versucht KOMA-Script, dies automatisch herauszufinden. (Da
% es sich diese Information in der .aux-Datei merkt, kann ein
% weiterer LATEX-Lauf notwendig sein.)
% Paket für Deutsche Sprache (Übersetzungen von Chapter zu Kapitel, 
% richtige Umlaute, richtige % Silbentrennung)
% siehe auch http://de.wikipedia.org/wiki/Babel-System
\usepackage[english,ngerman]{babel}

% Eingabecodierung, deutsche Umlaute oder die akzentuierten Zeichen sind verfügbar 
% und können direkt eingegeben werden
% siehe auch http://de.wikipedia.org/wiki/UTF8
\usepackage[utf8]{inputenc}

% Ausgabeschriftart von LaTeX festlegen
% siehe auch http://de.wikibooks.org/wiki/LaTeX-Schnellkurs:_Erste_Schritte
\usepackage[T2A,T1]{fontenc}
\DeclareUnicodeCharacter{0403}{{\fontencoding{T2A}\fontfamily{cmr}\selectfont\CYRGUP}}
\DeclareUnicodeCharacter{043D}{{\fontencoding{T2A}\fontfamily{cmr}\selectfont\cyrn}}

% schönere Schriftart
\renewcommand*{\ttdefault}{lmtt}
\usepackage{newpxtext}

% mikrotypografische Verbesserungen
\usepackage{microtype}
\microtypesetup{protrusion=true,expansion=true}
\DeclareMicrotypeSet{thesistext}{encoding={T1,TS1}}
\UseMicrotypeSet[expansion]{thesistext}

% Paket zum Erweitern der Tabelleneigenschaften
\usepackage{array}
\usepackage{longtable}

% Paket zum Überprüfen, ob ein pdf oder etwas anderes erstellt wird
\usepackage{ifpdf}

% Paket um Grafiken einbetten zu können
% erlaubt sind die Formate png, pdf, jpg
% nicht erlaubt sind eps/ps für pdftex/pdflatex
\usepackage{graphicx}
\usepackage{float}
% Standardpfad für Grafiken
\graphicspath{{bilder/}}
% Standardendung für Grafiken
\ifpdf
	% für die pdf-Ausgabe
	\DeclareGraphicsExtensions{.pdf,.jpg,.png}
\else
	% für die html-Ausgabe
	\DeclareGraphicsExtensions{.png,.jpg}
\fi

% ermöglicht Punkte in Dateinamen von Bildern
\usepackage{grffile}

\usepackage{xcolor}

% Paket zum "Zeichnen" von Grafiken (TikZ ist kein Zeichenprogramm)
\usepackage{tikz}
\usetikzlibrary{positioning,arrows.meta,fit,calc,shapes.misc,shapes.geometric}

% Paket zum Zeichnen von elektrischen Schaltkreisen mit TikZ
\usepackage[european,EFvoltages]{circuitikz}
\def\TikZ{Ti\emph{k}Z}
\def\Circuitikz{Circui\TikZ}

% Paket zum Erstellen von Plots mit TikZ
\usepackage{pgfplots}
% immer die neueste Version benutzen
\pgfplotsset{compat=newest}
% Verbindung von Linien durch eine schräge Kante
\pgfplotsset{every axis/.append style={line join=bevel}}
% neue Größenvorlage für zwei Plots nebeneinander anlegen
\pgfplotsset{
	scriptsize/.style={
		width=4.5cm,
		height=,
		legend style={font=\scriptsize},
		tick label style={font=\scriptsize},
		label style={font=\footnotesize},
		title style={font=\footnotesize},
		every axis title shift=0pt,
		max space between ticks=15,
		every mark/.append style={mark size=7},
		major tick length=0.1cm,
		minor tick length=0.066cm,
	},
}
% Legendeneintrage standardmäßig links ausrichten
\pgfplotsset{legend cell align=left}
% Anzahl der kleinen Teilstriche zwischen zwei großen Teilstrichen
%\pgfplotsset{minor x tick num={3}}
%\pgfplotsset{minor y tick num={3}}
% feines Gitternetz zeichnen
%\pgfplotsset{xminorgrids}
%\pgfplotsset{yminorgrids}
% nur nach den Achsen skalieren
\pgfplotsset{scale only axis}
% Farben wie in MATLAB definieren (Version R2014a und früher)
%%%%%%%% Farben wie in MATLAB definieren (ab Version R2014b)
% siehe: https://de.mathworks.com/help/matlab/graphics_transition/why-are-plot-lines-different-colors.html
% Die RGB-Werte können auch per
% >> figure;
% >> get(gca,'defaultAxesColorOrder');
% ausgelesen werden.
% Vielen Dank an Thomas Aab für die Aktualisierung.
% Farbreihenfolge von pgfplots benutzen
%\pgfplotsset{cycle list name=color list}
% nur Graustufen benutzen
%\pgfplotsset{cycle list name=linestyles}
% Strichstärke auf 1pt festlegen
\pgfplotsset{every axis plot/.append style={line width=1pt}}
% für deutsche Dokumente ein Komma benutzen
\addto\extrasngerman{\pgfplotsset{/pgf/number format/.cd,set decimal separator={{{,}}}}}
% ein halbes Leerzeichen als Tausendertrennzeichen benutzen
\pgfplotsset{/pgf/number format/.cd,1000 sep={\,}}

% Paket um mehrere Grafiken als eine Abbildung darstellen zu können
\usepackage{subfig}

% mathematische Symbole aus dem AMS (Amerikanische Mathematische Gesellschaft) Paket
\usepackage{amsmath}
\numberwithin{equation}{chapter}
% links-ausgerichtete Gleichungen
%\usepackage[fleqn]{amsmath}
\usepackage{newpxmath}
\usepackage{bm}
% kein zusätzlicher Abstand nach dem Dezimalkomma im Mathemodus
\usepackage{icomma}

% Paket um den Zeilenabstand zu ändern
\usepackage{setspace}
% Zeilenabstand 1,1 (wie die THB-LaTeX-Vorlage)
%\setstretch{1.1}
\onehalfspacing
\emergencystretch=2em
\clubpenalty=10000
\widowpenalty=10000
\displaywidowpenalty=10000

% Paket um Textteile drehen zu können
\usepackage{rotating}

% Paket für bessere Tabellen
\usepackage{booktabs}


% Seitenstil und Typografie werden in gestaltung.tex konfiguriert.
\usepackage[automark]{scrlayer-scrpage}

% Package für die deutsche BibTeX-Umgebung
% (sonst kann man keine Auflage eines Buches angeben)
\usepackage[backend=biber,style=numeric-comp,sorting=none,maxbibnames=99,maxcitenames=2,giveninits=true,
  uniquename=init,doi=true,isbn=false,url=true,eprint=true,date=year,urldate=long,urldateusetime=true,urltime=24h]{biblatex}

\DefineBibliographyStrings{ngerman}{
	andothers = {{et\addabbrvspace al\adddot}},
	urlfrom = {{Verfügbar unter}},
	urlseen = {Zugriff am},
}

% Datum und Uhrzeit des Abrufs mit Komma trennen (Leitfaden 4.5.3)
\DefineBibliographyExtras{ngerman}{\renewcommand*{\bibdatetimesep}{\addcomma\space}}

\DeclareFieldFormat*{url}{\bibstring{urlfrom}\addcolon\space\url{#1}}

\addbibresource{literatur.bib}  % Name deiner .bib-Datei anpassen
% Paket für bessere Anführungszeichen
\usepackage{csquotes}

% Paket für bessere Zitate


% Paket, um zu Verhindern, dass Zitate im Inhaltsverzeichnis, Abbildungsverzeichnis und Tabellenverzeichnis die Nummerierung beeinflussen
\usepackage{notoccite}

% Paket für Unterstreichungen
\usepackage[normalem]{ulem}

% Paket zur Änderung von Aufzählungen
\usepackage{enumitem}

% neue Umgebung für eine Aufzählung (unnummeriert) mit weniger vertikalen Abständen
\newenvironment{packed_itemize}{
\begin{itemize}[topsep=0pt,partopsep=0pt]
	\setlength{\itemsep}{0pt}
	\setlength{\parskip}{0pt}
	\setlength{\parsep}{0pt}
}{\end{itemize}}

% neue Umgebung für eine Aufzählung (nummeriert) mit weniger vertikalen Abständen
\newenvironment{packed_enumerate}{
\begin{enumerate}[topsep=0pt,partopsep=0pt]
	\setlength{\itemsep}{0pt}
	\setlength{\parskip}{0pt}
	\setlength{\parsep}{0pt}
}{\end{enumerate}}

% Paket für SI-Einheiten
\usepackage{siunitx}
% Optionen für das Paket
% Einstellungen in gestaltung.tex




% Paket zur Darstellung von Algorithmen
%\usepackage{algorithm2e}

% Paket für ein intelligenter Leerzeichen
\usepackage{xspace}

% ein paar Befehle für häufig benutzte Abkürzungen

% Markierung für Angaben, die der Verfasser noch bestätigen oder ergänzen muss
\newcommand{\pruefen}[1]{\textcolor{red!75!black}{[prüfen: #1]}}
\newcommand{\zB}{z.\,B.\xspace}
\newcommand{\dH}{d.\,h.\xspace} % den /dh-Befehl gibt es schon
% Kürzel der supervised Referenz (ÜBl. = überwachte Baseline). Nur hier ändern.
% Der Punkt entfällt, wenn im Text schon ein Punkt folgt (\UEBL. am Satzende ergibt ÜBl.).
\makeatletter
\DeclareRobustCommand{\UEBL}{ÜBl\@ifnextchar.{}{.\xspace}}
\makeatother

% neue mathematische Funktionen definieren
\DeclareMathOperator{\arcosh}{arcosh}
\DeclareMathOperator*{\argmax}{arg\,max}
\DeclareMathOperator*{\argmin}{arg\,min}

% Differentialoperator (kleines aufrechtes d)
\newcommand*{\diff}{\mathop{}\!\mathrm{d}}

% e-Funktion umdefinieren
\newcommand{\e}{\mathrm{e}}

% Indizes aufrecht (in Roman) schreiben
\newcommand{\ind}[1]{\mathrm{#1}}

% Punkt in einer Formelumgebung, z.B. am Ende einer Formel
\newcommand{\punkt}{\,\text{.}}
% ebenso Komma
\newcommand{\komma}{\,\text{,}}

% Betrag oder Absolutwert
\newcommand{\abs}[1]{\left| #1 \right|}

% Mittelwert oder Durchschnitt
\newcommand{\mean}[1]{\left\langle #1 \right\rangle}

% Paket für schönere Brüche im Textmodus
\usepackage{xfrac}
% Standardeinstellung wie im herkömmlichen Mathemodus mit schrägen Virgule-Bruchstrich
%\UseCollection{xfrac}{plainmath}

% Seitenränder: 3 cm auf allen Seiten (Kopfzeile liegt im oberen Rand, Seitenzahl der Kapitelanfänge im unteren)
\usepackage[margin=30mm,headheight=17pt,headsep=8mm,footskip=12mm]{geometry}
\raggedbottom

% bestimmte Trennung für bestimmte Wörter
\hyphenation{Mi-kro-strei-fen-lei-ters ver-schie-den-klas-si-ge ver-schie-den-klas-si-gen Me-cha-nis-mus}

% Paket für ein Symbolverzeichnis/Notationsverzeichnis
\usepackage[german]{nomencl}
\makenomenclature

% Paket für ein Abkürzungsverzeichnis
% - siehe http://de.wikibooks.org/wiki/LaTeX-W%C3%B6rterbuch:_Abk%C3%BCrzungsverzeichnis
% - ohne die nohyperlinks-Option werden zahlreiche Warnungen erzeugt
% - mit der nohyperlinks-Option werden aber keine Links mehr aus dem Text zum Abkürzungsverzeichnis erzeugt
\usepackage[nohyperlinks]{acronym}

% Abbildungen und Tabellen immer versuchen, 'hier' zu platzieren
\makeatletter
\renewcommand{\fps@figure}{htbp}
\renewcommand{\fps@table}{htbp}
\makeatother

% Legt den Anteil des Platzes am oberen Rand einer Seite fest, bis zu dem Gleitobjekte plaziert werden können.
% Standardwert: 0.7
\renewcommand{\topfraction}{0.85}

% Legt den Anteil des Platzes einer Seite fest, den Gleitobjekte (z.B. Abbildungen) am unteren Seitenrand einnehmen können. 
% Standardwert: 0.3
\renewcommand{\bottomfraction}{0.33}

% Einer dieser beiden Werte (\topfraction und \bottomfraction) sollte stets größer als \floatpagefraction sein!

% Definiert den Anteil einer Gleitobjekt-Seite, der belegt sein muß, bevor eine weitere Seite angelegt wird. 
% Standardwert: 0.5
\renewcommand{\floatpagefraction}{0.8} % eigene Seite nur für große Gleitobjekte, mittelgroße kommen oben auf eine Textseite

% Legt den Anteil des Platzes einer Seite mit Gleitobjekten fest, der mindestens für Text zur Verfügung steht. 
% Standardwert: 0.2
\renewcommand{\textfraction}{0.10}

% Gliederung bis subsection im Text
\setcounter{secnumdepth}{2}

% Gliederung bis subsection im Inhaltsverzeichnis
\setcounter{tocdepth}{2}
% Gestaltungsregeln getrennt vom Inhalt halten.

\setlength{\parskip}{0pt}
\setlength{\parindent}{2em}
\emergencystretch=1.5em
\clubpenalty=10000
\widowpenalty=10000
\displaywidowpenalty=10000
\raggedbottom
\setlength{\parfillskip}{0pt plus 0.75\textwidth}
\AtBeginBibliography{\setlength{\parfillskip}{0pt plus 1fil}}
\renewcommand*{\multicitedelim}{\addcomma}%

\setkomafont{disposition}{\normalfont\bfseries}
\setkomafont{chapter}{\LARGE}          %
\setkomafont{section}{\large}          %
\setkomafont{subsection}{\normalsize}  %
\setkomafont{subsubsection}{\normalsize}
\setkomafont{paragraph}{\normalsize}
\setkomafont{descriptionlabel}{\normalfont\bfseries}
\renewcommand*{\chapterformat}{\thechapter\enskip}
\renewcommand*{\sectionformat}{\thesection\enskip}
\renewcommand*{\subsectionformat}{\thesubsection\enskip}
\RedeclareSectionCommand[beforeskip=2\baselineskip,afterskip=2\baselineskip,afterindent=false]{chapter}
\RedeclareSectionCommand[beforeskip=1.5\baselineskip plus .3\baselineskip minus .2\baselineskip,afterskip=.75\baselineskip plus .1\baselineskip,afterindent=false]{section}
\RedeclareSectionCommand[beforeskip=1.25\baselineskip plus .2\baselineskip minus .2\baselineskip,afterskip=.5\baselineskip plus .1\baselineskip,afterindent=false]{subsection}
\RedeclareSectionCommand[beforeskip=1\baselineskip plus .2\baselineskip minus .2\baselineskip,afterskip=.5\baselineskip,afterindent=false]{subsubsection}
\RedeclareSectionCommand[beforeskip=.75\baselineskip plus .2\baselineskip,afterskip=-1em]{paragraph}

\clearpairofpagestyles
\automark[chapter]{chapter}
\renewcommand*{\chaptermarkformat}{\thechapter\enskip}
\ihead{\headmark}
\ohead{\pagemark}
\cfoot[\pagemark]{}
\setkomafont{pageheadfoot}{\normalfont\small}
\setkomafont{pagenumber}{\normalfont}
\setkomafont{pagehead}{\normalfont}
\setkomafont{headsepline}{\normalcolor}
\KOMAoptions{headsepline=0.4pt}

\DeclareTOCStyleEntry[beforeskip=.6\baselineskip plus 1pt,entryformat=\bfseries,pagenumberformat=\bfseries]{tocline}{chapter}
\DeclareTOCStyleEntries[beforeskip=0pt]{tocline}{section,subsection}
\AtBeginBibliography{\small\setstretch{1}\setlength{\bibitemsep}{.5\baselineskip}}
\AtEveryBibitem{\ifkeyword{freiefassung}{}{\iffieldundef{doi}{}{\clearfield{url}\clearfield{urlyear}\clearfield{urlmonth}\clearfield{urlday}}}}

\deffootnote{1.5em}{0pt}{\makebox[1.5em][l]{\thefootnotemark}}
\setkomafont{footnote}{\footnotesize}   %
\setlength{\footnotesep}{\baselineskip}

\captionsetup{format=plain,font=small,labelfont=bf,textfont=normalfont,
  labelsep=colon,justification=justified,singlelinecheck=false,
  margin=0pt,skip=6pt}
\captionsetup[figure]{position=bottom}
\captionsetup[table]{position=top}

\setlength{\floatsep}{12pt plus 3pt minus 2pt}
\setlength{\textfloatsep}{16pt plus 4pt minus 3pt}
\setlength{\intextsep}{14pt plus 3pt minus 2pt}
\makeatletter
\setlength{\@fptop}{0pt}
\setlength{\@fpsep}{\floatsep}
\setlength{\@fpbot}{0pt plus 1fil}
\makeatother

\renewcommand{\arraystretch}{1}
\setlength{\tabcolsep}{6pt}
\setlength{\heavyrulewidth}{0.08em}
\setlength{\lightrulewidth}{0.05em}
\setlength{\cmidrulewidth}{0.03em}
\AddToHook{env/table/begin}{\small\setstretch{1}}

\newenvironment{notationtable}{%
  \begingroup\small\setstretch{1}%
  \renewcommand{\arraystretch}{1.15}%
  \setlength{\LTpre}{6pt}\setlength{\LTpost}{6pt}%
  \begin{longtable}{@{}>{\raggedright\arraybackslash}p{0.29\textwidth}>{\raggedright\arraybackslash}p{\dimexpr0.71\textwidth-2\tabcolsep\relax}@{}}%
}{%
  \end{longtable}\endgroup
}
\newcommand*{\notationtablehead}{\textbf{Symbol} & \textbf{Bedeutung}\\}
\newcommand*{\acronymgroup}[1]{\textbf{#1}}

\setlist{leftmargin=*,topsep=.3\baselineskip,itemsep=.1\baselineskip,parsep=0pt}
\setlist[itemize,1]{label=\textbullet}
\setlist[itemize,2]{label=\textendash}
\setlist[enumerate,1]{label=\arabic*.}
\setlist[enumerate,2]{label=\alph*)}

\sisetup{locale=DE,group-minimum-digits=5,group-separator={\,},
  mode=math,reset-text-series=false,text-series-to-math=true,
  propagate-math-font=true,reset-math-version=false,
  uncertainty-mode=separate,per-mode=symbol,range-phrase={\ bis\ },
  table-align-text-before=false}

\definecolor{tolblue}{HTML}{4477AA}
\definecolor{tolred}{HTML}{EE6677}
\definecolor{tolgreen}{HTML}{228833}
\definecolor{tolpurple}{HTML}{AA3377}
\definecolor{tolyellow}{HTML}{CCBB44}
\definecolor{tolcyan}{HTML}{66CCEE}
\definecolor{tolgrey}{HTML}{BBBBBB}
\definecolor{neutral}{HTML}{4D4D4D}
\definecolor{thbrot}{HTML}{CE1628} %

\tikzset{
  diagram/.style={font=\footnotesize, line width=0.5pt, >={Latex[length=1.8mm,width=1.3mm]},
    every node/.append style={execute at begin node={\setstretch{1}}}},
  box/.style={draw=black, line width=0.5pt, rectangle, align=center, inner sep=3pt,
    minimum height=7mm, minimum width=18mm, fill=white},
  model/.style={box, fill=black!8},
  op/.style={draw=black, line width=0.5pt, circle, inner sep=1pt, minimum size=5mm, font=\footnotesize},
  arr/.style={->, line width=0.6pt, draw=black},
  aux/.style={->, line width=0.6pt, draw=black, dash pattern=on 2pt off 1.5pt},
  lbl/.style={font=\footnotesize, fill=white, inner sep=1pt},
  paneltitle/.style={font=\footnotesize, anchor=south west, inner sep=0pt},
  group/.style={draw=black, line width=0.4pt, dash pattern=on 2pt off 1.5pt, inner sep=4pt},
  pt/.style={circle, inner sep=0pt, minimum size=1.6mm, draw=black, line width=0.3pt},
  anchorpt/.style={pt, fill=black},
  pospt/.style={pt, fill=tolblue},
  negpt/.style={pt, fill=tolred},
  proxy/.style={rectangle, inner sep=0pt, minimum size=2mm, draw=black, line width=0.4pt, fill=white},
}

\pgfplotscreateplotcyclelist{thesis}{
  {neutral, mark=o},
  {tolblue, mark=square},
  {tolred, mark=triangle},
  {tolgreen, mark=diamond},
  {tolpurple, mark=pentagon},
  {tolyellow, mark=star}%
}
\pgfplotsset{
  every axis/.append style={font=\footnotesize, line width=0.5pt,
    tick style={line width=0.5pt, black}, tick align=inside,
    major tick length=3pt, legend style={draw=none, fill=none, font=\footnotesize},
    legend cell align=left, label style={font=\footnotesize},
    tick label style={font=\footnotesize}},
  every axis plot/.append style={line width=1pt, mark size=1.6pt},
  xmajorgrids=false, ymajorgrids=false,
  cycle list name=thesis,
}


% Hyperref-Paket für Links innerhalb des PDF Dokuments einbinden
% dieses Paket sollte als letztes eingebunden werden
% siehe auch http://www.math.uni-hamburg.de/home/iffland/Materialien/Einf_hyperref.pdf
\usepackage{xurl}
\urlstyle{same}
\Urlmuskip=0mu plus 2mu minus 0.2mu
\usepackage[
hidelinks,
pdftitle={Semi-Supervised Metric Learning mit Vision-Transformer-Modellen — Helfen ungelabelte Daten?},
pdfauthor={Sebastian Irmer},
pdfcreator={MiKTeX, LaTeX mit hyperref und KOMA-Script},
pdfsubject={Masterarbeit},
pdfkeywords={Metric Learning, Semi-Supervised Learning, Vision Transformer, Image Retrieval},
pdfstartview={Fit}
]{hyperref}
\usepackage{bookmark}

% Ende der Datei
